%%%%%%%%%%%%%%%%%%%%%%%%%%%%%%%%%%%%%%%%%%%%%%%%%%%%%%%%%%%%%%%%%%%%%%
% Lớp 12 | Các số đặc trưng ... ghép nhóm | Phương sai, độ lệch chuẩn của mẫu số liệu ghép nhóm
% Dạng 2D3?2-2: Tìm phương sai, độ lệch chuẩn
% Số câu: 43  (sinh tự động bởi tach_cau.py; từ nay thầy sửa trực tiếp file này)
%%%%%%%%%%%%%%%%%%%%%%%%%%%%%%%%%%%%%%%%%%%%%%%%%%%%%%%%%%%%%%%%%%%%%%

%%%%%%%%%%%%%%%% TRẮC NGHIỆM 4 PHƯƠNG ÁN %%%%%%%%%%%%%%%%

% nguồn: Lop_12_Dai_so_Chuong_3
\begin{ex}%[2D3N2-2]%[ma:1fea6179]
	Trong các khẳng định  sau, khẳng định nào sai?
	 \choice
		{Phương sai luôn  luôn là số không âm}
		{Phương sai là  bình phương của độ lệch chuẩn}
		{Phương sai càng lớn thì độ phân tác của các giá trị quanh số trung bình càng  lớn}
		{\True Phương sai  luôn luôn lớn hơn độ lệch chuẩn}
	\loigiai{ 
	Khi $ x\in\left(0;1\right)$ thì $s^2<s$ nên khẳng định  phương sai luôn lớn hơn độ lệch chuẩn là sai.}
\end{ex}

% nguồn: Lop_12_Dai_so_Chuong_3
\begin{ex}%[2D3N2-2]%[ma:e7835d24]
	Cho mẫu số liệu thống kê $\{2,4,6,8,10\}$. Phương sai của mẫu số liệu  trên là
	\choice 
		 {$40$}
		{$10$} 
		 {$6$}
		{\True  $8$}
	 \loigiai{Có $\bar{x}=\dfrac{1}{5}(2+4+6+8+10)=6$.\\
	 Phương sai là $s^2=\dfrac{1}{5}\left[(2-6)^2+(4-6)^2+(6-6)^2+(8-6)^2+(10-6)^2\right]=8$.}
\end{ex}

% nguồn: Lop_12_Dai_so_Chuong_3
\begin{ex}%[2D3H2-2]%[ma:b9237a6f]
	Một siêu thị thống kê số tiền (đơn vị: chục nghìn đồng) mà 44 khách hàng mua hàng ở siêu thị đó trong một  ngày. Số liệu được ghi lại trong bảng dưới đây.
	\begin{longtable}{|c|c|c|}
		 \hline
		Nhóm & Giá trị đại  diện & Tần số\\
		\hline 
		$[40;45)$ &  $42{,}5$ & $4$\\
		 \hline
		$[45;50)$ & $47{,}5$  & $14$\\
		 \hline
		 $[50;55)$ & $52{,}5$ & $8$\\
		\hline 
		$[55;60)$ & $57{,}5$  & $10$\\
		\hline 
		$[60;65)$ & $62{,}5$ &  $6$\\
		\hline 
		$[65;70)$ & $67{,}5$ &  $2$\\
		 \hline
		 & &$n=44$\\
		 \hline
	\end{longtable}
	Phương sai của mẫu số liệu ghép nhóm  trên là
	 \choice
		 {$53{,}2$}
		 {\True $46{,}1$}
		 {$30$}
		 {$11$}
	 \loigiai{
		Số trung bình cộng của mẫu số liệu ghép nhóm  là\\
		$\overline{x}=\dfrac{4\cdot 42{,}5+14\cdot 47{,}5+8\cdot 52{,}5+10\cdot 57{,}5+6\cdot 62{,}5+2\cdot  67{,}5}{44}=\dfrac{585}{11}$.\\
		Phương  sai của mẫu số liệu ghép nhóm là\\
		$\begin{aligned}
		&s^2=\dfrac{4\left(42{,}5-\dfrac{585}{11}\right)^2+14\left(47{,}5-\dfrac{585}{11}\right)^2+8\left(52{,}5-\dfrac{585}{11}\right)^2+10\left(57{,}5-\dfrac{585}{11}\right)^2}{44}\\ 
		 &+\dfrac{+6\left(62{,}5-\dfrac{585}{11}\right)^2+2\cdot \left(67{,}5-\dfrac{585}{11}\right)^2}{44}\approx 46{,}12.\\
\end{aligned}$}
\end{ex}

% nguồn: Lop_12_Dai_so_Chuong_3
\begin{ex}%[2D3H2-2]%[ma:d190bffb]
	Dũng là học sinh rất giỏi chơi rubik, bạn có thể giải nhiều loại khối rubik khác nhau. Trong một lần tập luyện giải khối  rubik $3\times 3$, bạn Dũng đã tự thống kê lại thời gian giải rubik trong $25$ lần giải liên tiếp ở bảng sau:
	\begin{center}
		\begin{tabular}{|c|c|c|c|c|c|}
			 \hline
			Thời gian giải rubik  (giây) &$[8;10)$ &$[10;12)$ &$[12;14)$ &$[14;16)$&$[16;18)$ \\
			 \hline
			Số  lần &$4$ &$6$ &$8$ &$4$&$3$ \\
			\hline 
		\end{tabular}
	\end{center}
	Độ lệch chuẩn của mẫu số liệu ghép nhóm có giá trị gần nhất  với giá trị nào dưới đây?
	 \choice
		 {$5{,}98$}
		 {$6$}
		{\True  $2{,}44$}
		{$2{,}5$} 
	 \loigiai{
		Ta có bảng sau 
		\begin{center}
			\begin{tabular}{|c|c|c|c|c|c|}
				\hline 
				Thời gian giải rubik (giây) &$[8;10)$ &$[10;12)$  &$[12;14)$ &$[14;16)$&$[16;18)$ \\
				\hline 
				Giá trị đại diện &$9$ &$11$ &$13$  &$15$&$17$ \\
				 \hline
				Số lần &$4$ &$6$ &$8$  &$4$&$3$ \\
				 \hline
			\end{tabular}
		\end{center}
		Số trung bình  của mẫu số liệu ghép nhóm là
		$$\overline{x}=\dfrac{4\cdot  9+6\cdot 11+8\cdot 13+4\cdot 15+3\cdot 17}{25}=12{,}68.$$
		Phương  sai của mẫu số liệu ghép nhóm là
		$$S^2=\dfrac{1}{25}\left[4\cdot 9^2+6\cdot 11^2+8\cdot 13^2+4\cdot 15^2+3\cdot 17^2\right]-(12{,}68)^2=5{,}9776.$$ 
		Độ lệch chuẩn của mẫu  số liệu ghép nhóm là
		$$S=\sqrt{S^2}=\sqrt{5{,}9776}\approx  2{,}44.$$
	 }
\end{ex}

% nguồn: Lop_12_Dai_so_Chuong_3
\begin{ex}%[2D3H2-2]%[ma:69d7b4be]
	Một siêu thị  thống kê số tiền (đơn vị: chục nghìn đồng) mà 44 khách hàng mua hàng ở siêu thị đó trong một ngày. Số liệu được ghi lại trong Bảng sau.
	\begin{center}
		\begin{tabular}{|c|c|c|}
			\hline 
			Nhóm & Giá  trị đại diện & Tần số \\
			 \hline
			$[40;45)$  & $42{,}5$ & $4$ \\
			\hline 
			$[45;50)$ & $47{,}5$  & $14$ \\
			\hline 
			$[50;55)$ &  $52{,}5$ & $8$ \\
			\hline 
			$[55;60)$ & $57{,}5$  & $10$ \\
			\hline 
			$[60;65)$ & $62{,}5$ & $6$ \\ 
			\hline 
			$[65;70)$ & $67{,}5$ &  $2$ \\
			 \hline
			 & & $n=44$ \\
			\hline 
		\end{tabular}
	\end{center}
	 Phương sai của mẫu số liệu ghép nhóm trên là (kết quả làm tròn đến hàng phần mười)
	 \choice
		{$53{,}2$} 
		 {$6{,}8$}
		{$47{,}2$} 
		{\True $46{,}1$} 
	\loigiai{ 
		Ta có $\bar{x}=\dfrac{42{,}5\cdot 4+47{,}5\cdot 14+52{,}5\cdot  8+57{,}5\cdot 10+62{,}5\cdot 6+67{,}5\cdot 2}{44}=\dfrac{585}{11}$.\\
		$s^2=\dfrac{(42{,}5-\bar{x})\cdot 4+(47{,}5-\bar{x})\cdot 14+(52{,}5-\bar{x})\cdot  8+(57{,}5-\bar{x})\cdot 10+(62{,}5-\bar{x})\cdot 6+(67{,}5-\bar{x})\cdot 2}{44}$\\
		$\approx  46{,}1$.
	 }
\end{ex}

% nguồn: Lop_12_Dai_so_Chuong_3
\begin{ex}%[2D3H2-2]%[ma:03ce88f3]
	Mỗi ngày bác Hương đều đi bộ để rèn luyện sức khoẻ. Quãng đường đi bộ mỗi ngày (đơn vị: km) của bác Hương trong $20$ ngày được thống kê lại ở bảng  sau
	\begin{longtable}{|c|c|c|c|c|c|}
		 \hline
		Quãng đường (km) & $[2{,}7;3{,}0)$ & $[3{,}0;3{,}3)$ & $[3{,}3;3{,}6)$  & $[3{,}6;3{,}9)$ & $[3{,}9;4{,}2)$\\
		 \hline
		Số ngày & $3$ & $6$ & $5$ & $4$ &  $2$\\
		 \hline
	\end{longtable}
	Độ  lệch chuẩn của mẫu số liệu ghép nhóm có giá trị gần nhất với giá trị nào dưới đây?
	\choice 
		{$3{,}41$} 
		{$11{,}62$} 
		{$0{,}017$} 
		{\True $0{,}36$} 
	 \loigiai{
		\begin{longtable}{|c|c|c|c|c|c|}
			\hline 
			Giá trị đại  diện & $2{,}85$ & $3{,}15$ & $3{,}45$ & $3{,}75$ & $4{,}05$\\
			 \hline
			Số  ngày & $3$ & $6$ & $5$ & $4$ & $2$\\
			\hline 
		\end{longtable}
		Số trung bình $\overline{x}=\dfrac{3\cdot 2{,}85+6\cdot 3{,}15+5\cdot  3{,}45+4\cdot 3{,}75+2\cdot 4{,}05}{20}=3{,}39$.\\
		Phương  sai $S^2=\dfrac{3\cdot 2{,}85^2+6\cdot 3{,}15^2+5\cdot 3{,}45^2+4\cdot 3{,}75^2+2\cdot 4{,}05^2}{20}-3{,}39^2=0{,}1314$.\\
		Độ lệch  chuẩn $s=\sqrt{0{,}1314}\approx 0{,}36$.
	} 
\end{ex}

% nguồn: Lop_12_Dai_so_Chuong_3
\begin{ex}%[2D3H2-2]%[ma:60434ba9]
	Bảng sau đây biểu diễn mẫu số  liệu ghép nhóm về thời gian sử dụng của một thiết bị điện tử (đơn vị: nghìn giờ).
	\begin{center}
		\begin{tabular}{|c|c|c|c|c|c|}
			\hline Nhóm & $[1{,}0; 1{,}2)$ &  $[1{,}2; 1{,}4)$ & $[1{,}4; 1{,}6)$ & $[1{,}6; 1{,}8)$ & $[1{,}8; 2{,}0)$ \\
			\hline Tần số & $20$ &  $34$ & $13$ & $10$ & $3$ \\
			 \hline
		\end{tabular}
	\end{center}
	Phương sai của  mẫu số liệu ghép nhóm trên thuộc khoảng
	\choice 
		 {$(0{,}04; 0{,}045)$}
		{\True  $(0{,}045; 0{,}05)$}
		 {$(0{,}05; 0{,}055)$}
		{$(0{,}055; 0{,}06)$} 
	\loigiai{ 
		 Số trung bình cộng của mẫu số liệu trên là
		\[\overline{x}=\dfrac{20\cdot 1{,}1+34\cdot 1{,}3+13\cdot  1{,}5+10\cdot 1{,}7+3\cdot 1{,}9}{80}=1{,}355.\]
		Phương sai của mẫu số liệu ghép nhóm trên là \begin{eqnarray*}
		s^2&=&\dfrac{20\left(1{,}1-1{,}355\right) ^2+34\left(1{,}3-1{,}355\right)  ^2+13\left(1{,}5-1{,}355\right) ^2}{80}\\
		&&+\dfrac{10\left(1{,}7-1{,}355\right) ^2+3\left(1{,}9-1{,}355\right) ^2  }{80}=0{,}046975\in (0{,}045; 0{,}05).
	\end{eqnarray*}
 }
\end{ex}

% nguồn: Lop_12_Dai_so_Chuong_3
\begin{ex}%[2D3H2-2]%[ma:77b3e7fd]
	Người ta ghi lại tiền lãi (đơn vị: triệu đồng) của một số nhà đầu tư (với số tiền đầu tư như nhau), khi đầu tư vào hai lĩnh vực $ A,B$ cho kết quả  như sau
	\begin{longtable}{|c|c|c|c|c|c|}
		\hline 
		Tiền lãi & $[5;10)$ &  $[10;15)$ & $[15;20)$ & $[20;25)$ & $[25;30)$\\
		 \hline
		Số nhà đầu tư vào lĩnh vực A & $2$ & $5$ & $8$ & $6$  & $4$\\
		\hline 
		Số nhà đầu tư vào lĩnh vực  B & $8$ & $4$ & $2$ & $5$ & $6$ \\
		\hline 
	\end{longtable}
	Tính  độ lệch chuẩn cho các mẫu số liệu về tiền lãi của các nhà đầu tư ở lĩnh vực B
	 \choice
		{\True $  8{,}039$}
		{$ 8{,}5314$} 
		{$ 8{,}6214$} 
		 {$ 8{,}5268$}
	\loigiai{ 
		Ta có mẫu số liệu ghép nhóm với giá tri đại diện  là
		\begin{longtable}{|c|c|c|c|c|c|}
			 \hline
			 Tiền lãi & $[5;10)$ & $[10;15)$ & $[15;20)$ & $[20;25)$ & $[25;30)$\\
			\hline 
			Giá trị đại diện & $7{,}5$ & $12{,}5$  & $17{,}5$ & $22{,}5$ & $27{,}5$\\
			 \hline
			Số nhà đầu tư vào lĩnh vực A &  $2$ & $5$ & $8$ & $6$ & $4$\\
			 \hline
			Số nhà đầu tư vào lĩnh vực B & $8$ & $4$ & $2$ & $5$  & $6$ \\
			 \hline
		\end{longtable}
		\noindent Phương sai của  mẫu số liệu về tiền lãi khi đầu tư vào lĩnh vực B:
		$$s_B^2=\dfrac{1}{25}\left(7{,}5^2\cdot 8+12{,}5^2\cdot 4+17{,}5^2\cdot 2+22{,}5^2\cdot  5+27{,}5^2\cdot 6\right)-16{,}9^2=64{,}64.$$
		Độ lệch chuẩn của mẫu số liệu về tiền lãi  khi đầu tư vào lĩnh vực B: $s_B=\sqrt{64{,}64}\approx 8{,}039$ (triệu đồng).
	 }
\end{ex}

% nguồn: Lop_12_Dai_so_Chuong_3
\begin{ex}%[2D3H2-2]%[ma:ef26251b]
	Dũng là học sinh rất giỏi chơi rubik, bạn có thể giải nhiều loại khối rubik khác nhau. Trong một lần tập luyện giải khối rubik $ 3\times 3$, bạn Dũng đã tự thống  kê lại thời gian giải rubik trong $25$ lần giải liên tiếp ở bảng sau
	\begin{longtable}{|c|c|c|c|c|c|}
		\hline 
		Thời gian giải rubik (giây) & $[8;10)$ & $[10;12)$ & $[12;14)$ & $[14;16)$  & $[16;18)$ \\
		 \hline
		Số  lần & $4$ & $6$ & $8$ & $4$ & $3$\\
		\hline 
	\end{longtable}
	Độ lệch chuẩn của mẫu số liệu ghép nhóm có giá trị gần nhất với  giá trị nào dưới đây?
	\choice 
		{\True $5{,}98$} 
		{$6$} 
		{$2{,}44$} 
		{$2{,}5$} 
	\loigiai{ 
		\begin{longtable}{|c|c|c|c|c|c|}
			\hline 
			Giá trị đại diện & $9$ & $11$ &  $13$ & $15$ & $17$ \\
			\hline 
			Số lần & $4$ & $6$ & $8$ & $4$  & $3$\\
			 \hline
		\end{longtable}
		Số trung bình $\overline{x}=\dfrac{4\cdot  9+6\cdot 11+8\cdot 13+4\cdot 15+3\cdot 17}{25}=12{,}68$.\\
	Độ lệch chuẩn $s=\sqrt{\dfrac{4\cdot 9^2+6\cdot  11^2+8\cdot 13^2+4\cdot 15^2+3\cdot 17^2}{25}-12{,}68^2}\approx 5{,}98$.}
\end{ex}

% nguồn: Lop_12_Dai_so_Chuong_3
\begin{ex}%[2D3H2-2]%[ma:ce4116a4]
	Khối lượng của một số quả trứng gà trong một trang trại  được ghi lại ở bảng sau:
	\begin{center}
		\centerline{ 
		\begin{tabular}{|c|c|c|c|c|}
			 \hline
			Khối lượng (gam) & $[43;45)$ & $[45;47)$ & $[47;49)$ &  $[49;51)$\\
			\hline 
			Số quả trứng & $12$ & $34$  & $32$ & $1$\\
			\hline 
		\end{tabular}}
	\end{center}
	Độ lệch chuẩn của mẫu số liệu ghép nhóm trên thuộc  nửa khoảng nào dưới đây?
	 \choice
		{\True $[1{,}4;1{,}6)$} 
		{$[2{,}0;2{,}2)$} 
		 {$[3{,}3;3{,}5)$}
		 {$[3{,}5;3{,}7)$}
	\loigiai{ 
		Ta có $n=12+34+32+1=79$.\\ 
		Giá trị đại diện cho các nhóm lần lượt là  $x_1=44$, $x_2=46$, $x_3=48$, $x_4=50$.\\
		Số trung bình của mẫu số liệu ghép nhóm  là
		$$\overline{x}=\dfrac{12 \cdot 44+34 \cdot 46  + 32 \cdot 48 + 1\cdot 50}{79}=\dfrac{3678}{79}.$$
		Phương sai của  mẫu số liệu ghép nhóm là\\
		$s^2=\dfrac{12\left( 44-\dfrac{3678}{79}\right) ^2+34\left(  46-\dfrac{3678}{79}\right) ^2+32\left( 48-\dfrac{3678}{79}\right) ^2+1\left( 50-\dfrac{3678}{79}\right) ^2}{79}=\dfrac{13232}{6241}$.\\
		Độ lệch chuẩn của mẫu số liệu  ghép nhóm là $s=\sqrt{\dfrac{13232}{6241}} \approx 1{,}456$.
	 }
\end{ex}

% nguồn: Lop_12_Dai_so_Chuong_3
\begin{ex}%[2D3H2-2]%[ma:58ce44af]
	Bảng sau thống kê khối lượng một số quả táo được lựa chọn  ngẫu nhiên từ một lô hàng (đơn vị: gam).
	\begin{center}
		\begin{tabular}{|c|c|c|c|c|c|}
			\hline Khối lương & $[80; 82)$ & $[82; 84)$ & $[84; 86)$ & $[86; 88)$ & $[88;  90)$ \\
			\hline Tần số & $18$ & $20$ & $24$  & $15$ & $13$ \\
			\hline 
		\end{tabular}
	\end{center}
	Phương sai của mẫu số liệu ghép nhóm trên (làm tròn đến hàng  phần trăm) là
	 \choice
		 {$2{,}64$}
		{\True  $6{,}96$}
		{$84{,}92$} 
		{$10$} 
	\loigiai{ 
		Số trung bình của  mẫu số liệu ghép nhóm trên (làm tròn đến hàng phần trăm) là
		\[\overline{x}=\dfrac{18\cdot 81+20\cdot  83+24\cdot 85+15\cdot 87+13\cdot 89}{90}=\dfrac{254}{3}\approx 84{,}67.\]
		Phương sai của mẫu số liệu ghép nhóm trên (làm tròn đến hàng phần  trăm) là
		\begin{eqnarray*}
			s^2&=&\dfrac{18\left( 81-\dfrac{254}{3}\right)^2 +20\left( 83-\dfrac{254}{3}\right) ^2+24\left( 85-\dfrac{254}{3}\right) ^2}{90}\\ 
			&&+\dfrac{15  \left( 87-\dfrac{254}{3}\right) ^2+13\left( 89-\dfrac{254}{3}\right) ^2}{90}
			=\dfrac{313}{45}\approx  6{,}96.
		\end{eqnarray*}
	 }
\end{ex}

% nguồn: Lop_12_Dai_so_Chuong_3
\begin{ex}%[2D3H2-2]%[ma:b17cc060]
	Để đánh giá chất lượng của một loại pin điện  thoại mới, người ta ghi lại thời gian nghe nhạc liên tục của điện thoại được sạc đầy pin cho đến khi hết pin cho kết quả sau
	\begin{center}
		\begin{tabular}{|l|c|c|c|c|c|}
			\hline Thời gian (giờ) &{$[5; 5{,}5)$}&{$[5{,}5; 6)$}&{$[6; 6{,}5)$}&{$[6{,}5; 7)$}&{$[7; 7{,}5)$}\\ 
			\hline Số chiếc điện thoại (tần  số) & $2$ & $8$ & $15$ & $10$ & $5$ \\
			\hline 
		\end{tabular}
	\end{center}
	Tính độ lệch chuẩn của mẫu số liệu ghép nhóm trên (làm tròn đến $4$ chữ số thập phân) 
	\choice 
		 {$0{,}4252$}
		{$0{,}5314$} 
		{$0{,}6214$} 
		{\True  $0{,}5268$}
	 \loigiai{
		\begin{center}
			\begin{tabular}{|l|c|c|c|c|c|}
				\hline Thời gian (giờ) & {$[5; 5{,}5)$} & {$[5{,}5; 6)$} & {$[6; 6{,}5)$}  & {$[6{,}5; 7)$} & {$[7; 7{,}5)$} \\
				\hline Giá  trị đại diện & $5{,}25$ & $5{,}75$ & $6{,}25$ & $6{,}75$ & $7{,}25$ \\
				\hline Số  chiếc điện thoại (tần số) & $2$ & $8$ & $15$ & $10$ & $5$ \\
				\hline 
			\end{tabular}
		\end{center}
		Số trung bình của  mẫu số liệu
		 $$
		\overline{x}=\dfrac{m_1 \cdot x_1+\ldots+m_k \cdot x_k}{n}=\dfrac{2\cdot5{,}25+8\cdot5{,}75+15\cdot6{,}25+10\cdot6{,}75+5\cdot7{,}25}{40}=6{,}35. 
		$$ 
		Phương sai của mẫu  số liệu ghép nhóm
		 $$
		s^2=\dfrac{1}{40}\left(2\cdot5{,}25^2+8\cdot5{,}75^2+15\cdot6{,}25^2+10\cdot6{,}75^2+5\cdot7{,}25^2\right)-6{,}35^2=0{,}2775. 
		$$ 
		Độ lệch chuẩn của mẫu  số liệu ghép nhóm là $s=\sqrt{s^2}=\sqrt{0{,}2775} \approx 0{,}5268$.
	 }
\end{ex}

% nguồn: Lop_12_Dai_so_Chuong_3
\begin{ex}%[2D3H2-2]%[ma:a00d14e6]
	Để đánh giá chất lượng của một loại pin điện thoại mới, người ta ghi lại thời gian nghe nhạc liên tục của điện thoại được sạc đầy pin cho đến khi hết pin cho kết  quả sau
	\begin{longtable}{|c|c|c|c|c|c|}
		\hline 
		Thời gian (giờ) & $[5;5{,}5)$ & $[5{,}5;6)$ & $[6;6{,}5)$  & $[6{,}5;7)$ & $[7;7{,}5)$\\
		\hline 
		Số chiếc điện thoại (tần số) & $2$ & $8$ & $15$ & $10$  & $5$\\
		 \hline
	\end{longtable}
	Tính độ lệch chuẩn của mẫu số liệu ghép  nhóm trên (làm tròn đến 4 chữ số thập phân)
	\choice 
		{\True $ 0{,}4252$} 
		{$  0{,}5314$}
		{$ 0{,}6214$} 
		{$  0{,}5268$}
	\loigiai{ 
		\begin{longtable}{|c|c|c|c|c|c|}
			 \hline
			Thời gian (giờ) & $[5;5{,}5)$ & $[5{,}5;6)$ & $[6;6{,}5)$ & $[6{,}5;7)$ &  $[7;7{,}5)$\\
			 \hline
			Giá trị đại diện & $5{,}25$ & $5{,}75$ & $6{,}25$ &  $6{,}75$ & $7{,}25$\\
			 \hline
			Số chiếc điện thoại (tần số) & $2$  & $8$ & $15$ & $10$ & $5$\\
			\hline 
		\end{longtable}
		\noindent Số  trung bình của mẫu số liệu
		$$\overline{x}=\dfrac{m_1\cdot x_1+\ldots +m_k\cdot x_k}{n}=\dfrac{2\cdot  5{,}25+8\cdot 5{,}75+15\cdot 6{,}25++10\cdot 6{,}75+5\cdot 7{,}25}{40}=6{,}35.$$
		Phương sai của mẫu số liệu ghép nhóm:\\ 
		$s^2=\dfrac{1}{40}\left(2\cdot 5{,}25^2+8\cdot 5{,}75^2+15\cdot  6{,}25^2+10\cdot 6{,}75^2+5\cdot 7{,}25^2\right)-6{,}35^2=0{,}2775$.\\
		Độ lệch chuẩn của  mẫu số liệu ghép nhóm là $ s=\sqrt{s^2}=\sqrt{0{,}2775}\approx 0{,}5268$.
	} 
\end{ex}

% nguồn: Lop_12_Dai_so_Chuong_3
\begin{ex}%[2D3H2-2]%[ma:d342f689]
	\immini{ 
	Xét mẫu số liệu ghép nhóm được  cho ở bảng 4. Phương sai của mẫu số liệu ghép nhóm đó là
	 }{
	\begin{tabular}{|c|c|c|}
		\hline 
		 Nhóm & Tần số & Tần số tích lũy\\
		 \hline
		 $[50;60)$ &$3$&$3$\\
		\hline 
		$[60;70)$  &$5$&$8$\\
		 \hline
		$[70;80)$  &$25$&$33$\\
		 \hline
		$[80;90)$ &$4$&$37$\\ 
		 \hline
		$[90;100)$  &$3$&$40$\\
		 \hline
		& $n=40$ &\\ 
		 \hline
	\end{tabular}
	 }
	 \choice
		{$9{,}08$} 
		 {\True $82{,}4375$}
		{$74{,}75$} 
		{$50$} 
	 \loigiai{
		Ta  có
		\begin{center}
			\resizebox{0.8\linewidth}{!}{\begin{tabular}{|c|c|c|c|c|c|}
			 \hline
			Nhóm  & $[50;60)$& $[60;70)$ &$[70;80)$& $[80;90)$&$[90;100)$\\
			 \hline
			Giá trị  đại diện & $55$ &$65$&$75$&$85$&$95$\\
			\hline 
			Tần  số & $3$&$5$&$25$&$4$&$3$\\
			\hline 
		\end{tabular}}
	\end{center}
	Số trung  bình của mẫu số liệu ghép nhóm là
	 \[
	\overline{x}=\dfrac{1}{40}\left(3\cdot 55 +5\cdot 65+25\cdot 75+4\cdot  85+3\cdot 95\right) = 74{,}75.
	\] 
	Phương sai của mẫu số liệu ghép  nhóm
	\begin{eqnarray*}
		s^2&=&  \dfrac{1}{40}\big[3\cdot(55-74{,}75)^2+5\cdot(65-74{,}75)^2+25\cdot(75-74{,}75)^2\\
		&&+\,  4\cdot(85-74{,}75)^2+3\cdot(95-74{,}75)^2\big]\\
		 &=&\dfrac{1\,319}{16}=82{,}4375.
	\end{eqnarray*}
} 
\end{ex}

% nguồn: Lop_12_Dai_so_Chuong_3
\begin{ex}%[2D3H2-2]%[ma:65e46381]
	Mỗi ngày bác Hương đều đi bộ để rèn luyện sức khoẻ. Quãng đường đi bộ mỗi ngày (đơn vị:  km) của bác Hương trong $20$ ngày được thống kê lại ở bảng sau
	\begin{longtable}{|c|c|c|c|c|c|}
		\hline 
		Quãng đường (km) & $[2{,}7;3{,}0)$ & $[3{,}0;3{,}3)$ & $[3{,}3;3{,}6)$ & $[3{,}6;3{,}9)$  & $[3{,}9;4{,}2)$\\
		 \hline
		Số ngày & $3$ & $6$ & $5$ & $4$  & $2$\\
		 \hline
	\end{longtable}
	Phương sai  của mẫu số liệu ghép nhóm là
	\choice 
		 {$3{,}39$}
		{$11{,}62$} 
		{\True $0{,}1314$} 
		{$0{,}36$} 
	 \loigiai{
		\begin{longtable}{|c|c|c|c|c|c|}
			 \hline
			Giá trị đại diện & $2{,}85$ & $3{,}15$ & $3{,}45$ & $3{,}75$ & $4{,}05$\\ 
			 \hline
			Số ngày & $3$ & $6$ & $5$ & $4$  & $2$\\
			\hline 
		\end{longtable}
		Số trung  bình $\overline{x}=\dfrac{3\cdot 2{,}85+6\cdot 3{,}15+5\cdot 3{,}45+4\cdot 3{,}75+2\cdot 4{,}05}{20}=3{,}39$.\\
		 Phương sai $S^2=\dfrac{3\cdot 2{,}85^2+6\cdot 3{,}15^2+5\cdot 3{,}45^2+4\cdot 3{,}75^2+2\cdot 4{,}05^2}{20}-3{,}39^2=0{,}1314$.
	 }
\end{ex}

% nguồn: Lop_12_Dai_so_Chuong_3
\begin{ex}%[2D3H2-2]%[ma:56a09feb]
	Một siêu thị thống kê số tiền (đơn vị: chục nghìn đồng) mà $44$ khách hàng mua  hàng ở siêu thị đó trong một ngày. Số liệu được ghi lại trong bảng dưới đây.
	\begin{longtable}{|c|c|c|}
		 \hline
		Nhóm  & Giá trị đại diện & Tần số\\
		 \hline
		$[40;45)$ & $42{,}5$ & $4$\\ 
		 \hline
		$[45;50)$ &  $47{,}5$ & $14$\\
		 \hline
		 $[50;55)$ & $52{,}5$ & $8$\\
		 \hline
		$[55;60)$ & $57{,}5$ & $10$\\ 
		\hline 
		$[60;65)$ &  $62{,}5$ & $6$\\
		 \hline
		 $[65;70)$ & $67{,}5$ & $2$\\
		\hline 
		&  &$n=44$\\
		\hline 
	\end{longtable}
	Độ lệch chuẩn của mẫu số liệu ghép nhóm trên (làm tròn  kết quả đến hàng đơn vị) là
	\choice 
		{\True $6{,}8$} 
		 {$7{,}3$}
		{$3{,}3$} 
		{$46{,}1$} 
	\loigiai{ 
		Số trung  bình cộng của mẫu số liệu ghép nhóm là\\
		$\overline{x}=\dfrac{4\cdot 42{,}5+14\cdot 47{,}5+8\cdot 52{,}5+10\cdot 57{,}5+6\cdot 62{,}5+2\cdot 67{,}5}{44}=\dfrac{585}{11}$.\\ 
		 Phương sai của mẫu số liệu ghép nhóm là\\
		$\begin{aligned}
		&s^2=\dfrac{4\left(42{,}5-\dfrac{585}{11}\right)^2+14\left(47{,}5-\dfrac{585}{11}\right)^2+8\left(52{,}5-\dfrac{585}{11}\right)^2+10\left(57{,}5-\dfrac{585}{11}\right)^2}{44}\\ 
		&+\dfrac{6\left(62{,}5-\dfrac{585}{11}\right)^2+2\cdot \left(67{,}5-\dfrac{585}{11}\right)^2}{44}\approx 46{,}12.\\ 
	\end{aligned}$\\
	Độ lệch chuẩn của  mẫu số liệu ghép nhóm là $ s=\sqrt{s^2}=\sqrt{46{,}12}\approx 6{,}8$.
} 
\end{ex}

% nguồn: Lop_12_Dai_so_Chuong_3
\begin{ex}%[2D3H2-2]%[ma:039e66cf]
	Bảng dưới đây thống kê độ ẩm không khí trung bình mỗi ngày đo được tại một trạm quan trắc trong 60  ngày đầu năm 2023.
	\begin{center}
		\resizebox{0.8\linewidth}{!}{\begin{tikzpicture}
		\draw[opacity=.25,thin,step=.2,cyan](0,0)  grid (7,2.5);
		\draw[opacity=1,cyan](0,0) grid (7,2.5); 
		\draw[stealth-stealth](0,2.5) node[left]{Tần số tương đối  (\%)}|-(7,0) node[below,anchor=west]{Độ ẩm (\%)};
		\foreach\x/\dientich[count=\i from 1] in {0/60,1/64,1/68,1.5/72,1/76,0.5/80}{ 
			\draw[pattern=north west lines,pattern color=cyan!65](\i-1,0)  rectangle +(1,\x);
			 \draw (\i,0) node[below]{\dientich};
		 }
		\foreach \y [count=\i from  1] in {10,20,30}{
			 \draw (-.1,\i/2)--(.1,\i/2)
			(0,\i/2)  node[left]{$\y$};}
		 \foreach \z [count=\i from 1] in {20,20,30,20,10}{
			 \draw (\i+0.5,\z/20) node[above] {$\z$};}
	\end{tikzpicture}
		 }
	\end{center}
	Độ lệch chuẩn  của mẫu số liệu ghép nhóm trên (làm tròn đến hàng phần nghìn) là
	 \choice
		{\True $4{,}996$} 
		 {$24{,}960$}
		{$5{,}038$} 
		{$25{,}383$} 
	\loigiai{Ta có bảng thống kê theo giá trị  đại diện như sau
		\begin{center}
			\begin{tabular}{|c|c|c|c|c|c|c|}
				 \hline
				Gía trị đại diện&  $62$ & $66$ & $70$ & $74$ & $78$ \\
				\hline 
				Tần số &$12$& $12$ & $18$ &$12$ & $6$  \\
				\hline 
			\end{tabular}
		\end{center}
		Cỡ mẫu là $n=60$. Số trung bình của mẫu số  liệu ghép nhóm trên là \[\overline{x}=\dfrac{12\cdot 62+12\cdot 66+18\cdot 70+12\cdot 74+6\cdot 78}{60}=69{,}2.\]
		Phương sai của mẫu số liệu ghép nhóm là 
		\[S^2=\dfrac{1}{60}\left(12\cdot 62^2+12\cdot 66^2+18\cdot 70^2+12\cdot 74^2+6\cdot 78^2\right)-69{,}2^2  =24{,}96.\]
		Độ lệch chuẩn của mẫu số liệu  ghép nhóm là
	\[S =  \sqrt{24{,}96} \approx 4{,}996.\]}
\end{ex}

% nguồn: Lop_12_Dai_so_Chuong_3
\begin{ex}%[2D3H2-2]%[ma:e1ec7dad]
	Bạn Chi rất thích nhảy hiện đại. Thời gian tập nhảy  mỗi ngày trong thời gian gần đây của bạn Chi được thống kê lại ở bảng sau
	\begin{longtable}{|c|c|c|c|c|c|}
		\hline 
		Thời gian (phút) & $[20;25)$ & $[25;30)$ & $[30;35)$ &  $[35;40)$ & $[40;45)$\\
		 \hline
		Số ngày & $6$ & $6$  & $4$ & $1$ & $1$ \\
		\hline 
	\end{longtable}
	Phương sai của mẫu số liệu ghép nhóm có giá trị gần  nhất với giá trị nào dưới đây?
	\choice 
		{$31{,}77$} 
		 {$32$}
		 {$31$}
		 {\True $31{,}44$}
	\loigiai{ 
		Số trung bình $\overline{x}=\dfrac{6\cdot 22{,}5+6\cdot 27{,}5+4\cdot  32{,}5+37{,}5+42{,}5}{18}\approx 28{,}33$.\\
		Phương sai $S^2=\dfrac{6\cdot 22{,}5^2+6\cdot  27{,}5^2+4\cdot 32{,}5^2+37{,}5^2+42{,}5^2}{18}-28{,}33^2=31{,}25$.
	 }
\end{ex}

% nguồn: Lop_12_Dai_so_Chuong_3
\begin{ex}%[2D3H2-2]%[ma:ad665cda]
	Điểm kiểm tra $15$ phút của $36$ học sinh lớp  11A được cho bởi bảng tần số ghép nhóm sau
	\begin{center}
		\begin{tabular}{|c|c|}
			\hline 
			 \textbf{Nhóm điểm} & \textbf{Tần số} \\
			 \hline
			$[1;3)$ & $3$  \\
			 \hline
			$[3;5)$ &  $2$ \\
			\hline 
			$[5;7)$ & $10$  \\
			\hline 
			$[7;9)$ &  $14$ \\
			 \hline
			 $[9;11)$ & $7$ \\
			 \hline
			& $n=36$  \\
			 \hline
		\end{tabular}
	\end{center}
	Độ lệch chuẩn bảng  ghép lớp trên là giá trị nào sau đây?
	 \choice
		{\True $2{,}23$} 
		{$3{,}61$} 
		 {$4{,}29$}
		{$3{,}85$} 
	\loigiai{ 
		\begin{center}
			\begin{tabular}{|c|c|c|}
				\hline 
				\textbf{Nhóm điểm} & \textbf{Tần số} & \textbf{Giá trị đại diện} \\ 
				 \hline
				$[1;3)$ &  $3$ & $2$\\
				 \hline
				$[3;5)$ & $2$ &  $4$\\
				\hline 
				$[5;7)$ & $10$ &  $6$\\
				\hline 
				$[7;9)$ & $14$  & $8$\\
				 \hline
				$[9;11)$ & $7$ &  $10$\\
				 \hline
				& $n=36$  \\
				\hline 
			\end{tabular}
		\end{center}
		Ta có số trung  bình $\overline{x}=\dfrac{3\cdot2+2\cdot4+10\cdot6+14\cdot8+7\cdot10}{36}=\dfrac{64}{9}$.\\
		Khi đó 
		\begin{align*}
			 s&=\sqrt{\dfrac{n_1\left(x_1-\overline{x}\right)^2+n_2\left(x_2-\overline{x}\right)^2+\cdots+n_m{\left(x_m-\overline{x}\right)^2}}{n}}\\
			&=\sqrt{\dfrac{3\left(2-\dfrac{64}{9}\right)^2+2\left(4-\dfrac{64}{9}\right)^2+10\left(6-\dfrac{64}{9}\right)^2+14\left(8-\dfrac{64}{9}\right)^2+7\left(10-\dfrac{64}{9}\right)^2}{36}}\\ 
			&\approx  2{,}23.
		\end{align*}
	 }
\end{ex}

% nguồn: Lop_12_Dai_so_Chuong_3
\begin{ex}%[2D3H2-2]%[ma:f6d867f7]
	Khảo sát thời gian sử dụng điện thoại $1$ ngày của học sinh lớp $12$A$1$ thu được mẫu số liệu ghép nhóm  sau
	\begin{center}
		\begin{tabular}{|c|c|c|c|c|c|}
			 \hline
			Thời gian (phút) & $[0; 20)$ & $[20; 40)$ & $[40; 60)$ & $[60; 80)$ & $[80;  100)$ \\ \hline
			Số học sinh & $2$ & $5$ &  $7$ & $19$ & $9$ \\ \hline
		\end{tabular}
	\end{center}
	Độ lệch chuẩn  của mẫu số liệu ghép nhóm trên thuộc khoảng nào dưới đây?
	 \choice
		{$(18; 20)$} 
		 {$(22; 24)$}
		{\True $(20;  22)$}
		{$(24; 26)$} 
	\loigiai{Cỡ mẫu $n=2+5+7+19+9=42$.\\ 
		Ta có bảng 
		\begin{center}
			\begin{tabular}{|c|c|c|c|c|c|}
				\hline 
				Thời  gian (phút) & $[0; 20)$ & $[20; 40)$ & $[40; 60)$ & $[60; 80)$ & $[80; 100)$ \\ \hline
				Giá trị đại diện  & $10$ & $30$ & $50$ & $70$ & $90$ \\ \hline
				Số học sinh &  $2$ & $5$ & $7$ & $19$ & $9$ \\ \hline
			\end{tabular}
		\end{center}
		Số trung bình cộng  là
		\[\overline{x}=\dfrac{10 \cdot 2 + 30 \cdot 5 + 50 \cdot 7 + 70 \cdot 19 + 90 \cdot 9}{42} \approx  63{,}3.\]
		Phương  sai là
		 \[s^2 = \dfrac{2 \cdot (10-63{,}3)^2+5 \cdot (30-63{,}3)^2+7 \cdot (50-63{,}3)^2+19 \cdot (70-63{,}3)^2+9 \cdot (90-63{,}3)^2}{42} \approx 469{,}8.\]
		Độ lệch chuẩn là $s = \sqrt{469{,}8} \approx 21{,}7 \in (22;  24)$.
	} 
\end{ex}

% nguồn: Lop_12_Dai_so_Chuong_3
\begin{ex}%[2D3V2-2]%[ma:71b43ec9]
	 \immini{
	Bảng bên biểu diễn mẫu số liệu ghép nhóm về chi tiêu bình quân (đơn vị: USD) của một lượt khách quốc tế đến Việt Nam phân theo $27$  quốc tịch năm $2019$.
	\begin{flushright}
		 (\textit{Nguồn: https://www.gso.gov.vn})
	\end{flushright}
	Độ lệch chuẩn của mẫu số  liệu đó nằm trong khoảng nào dưới đây?
	\choice[2] 
		 {$(200;300)$}
		{\True $(300;400)$} 
		 {$(400;500)$}
		{$(500;600)$} 
		 }{
		\begin{tabular}{|l|c|}
			\hline 
			Nhóm &  Tần số \\
			 \hline
			$[0; 500)$ &  $1$ \\
			 $[500; 1\,000)$ & $9$ \\
			$[1\,000; 1\,500)$  & $14$ \\
			$[1\,500; 2\,000)$ &  $2$ \\
			$[2\,000; 2\,500)$ & $1$ \\ 
			 \hline
	\end{tabular}
	 }
	\loigiai{ 
		Ta  có bảng sau
		\begin{center}
			\begin{tabular}{|l|c|c|}
				\hline 
				Nhóm &Giá trị đại diện  &Tần số \\
				\hline 
				$[0; 500)$& $250$  & $1$ \\
				$[500; 1\,000)$& $750$  & $9$ \\
				$[1\,000; 1\ 500)$& $1\,250$ &  $14$ \\
				$[1\,500; 2\,000)$& $1\,750$ & $2$ \\ 
				$[2\,000;  2\,500)$& $2\,250$ & $1$ \\
				\hline 
			\end{tabular}
		\end{center}
		Ta có giá trị  trung bình của mẫu số liệu là
		 \allowdisplaybreaks
		\begin{eqnarray*}
			\overline{x}=\dfrac{250 \cdot 1 + 750 \cdot 9 + 1\,250 \cdot 14 + 1\,750 \cdot 2  + 2\,250 \cdot 1}{27}= \dfrac{30\,250}{27}.
		\end{eqnarray*}
		Suy ra độ lệch chuẩn  của mẫu số liệu là \allowdisplaybreaks
		\begin{eqnarray*}
			&s &=\sqrt{\dfrac{1\cdot (250 - \overline{x})^2+ 9 \cdot (750 - \overline{x})^2 + 14\cdot (1\, 250 - \overline{x})^2 +  2 \cdot (1\, 750 - \overline{x})^2 + 1 \cdot (2\, 250 -\overline{x})^2}{27}}\\
			&&\approx  398{,}9.
		\end{eqnarray*}
	 }
\end{ex}


%%%%%%%%%%%%%%%% ĐÚNG/SAI %%%%%%%%%%%%%%%%

% nguồn: Lop_12_Dai_so_Chuong_3
\begin{ex}%[2D3H2-2]%[ma:75ac90f5]
	\immini{Kết quả kiểm tra cân nặng của $20$ học  sinh nam lớp 12A (làm tròn kết quả đến hàng đơn vị của kilôgam) được cho bởi bảng bên
	\choiceTF 
		{\True Khoảng biến thiên của  mẫu số liệu ghép nhóm đã cho là $20$}
		{\True Số trung bình cộng của mẫu số liệu được tính bởi  công thức\\
		$\overline{x}=\dfrac{8 \cdot 62+ 9 \cdot 66+1 \cdot  70+1 \cdot 74+ 1\cdot 78}{20}$}
		{Phương sai của mẫu số liệu  ghép nhóm đã cho là $s^2=\sqrt{\dfrac{436}{25}}$}
		{Độ lệch chuẩn của mẫu số liệu ghép nhóm đã cho (làm tròn kết quả đến hàng phần mười của kilôgam) là $4$ kg}}{\begin{tabular}{|>{\centering\arraybackslash}p{1.5cm}|>{\centering\arraybackslash}p{2cm}|c|}
		\hline 
		\textbf{Nhóm}&\textbf{Giá trị  đại diện}&\textbf{Tần số}\\
		\hline 
		$[60;64)$&$62$&$8$\\ 
		 \hline
		$[64;68)$&$66$&$9$  \\
		 \hline
		 $[68;72)$&$70$&$1$ \\
		\hline 
		$[72;76)$&$74$&$1$  \\
		 \hline
		$[76;80)$&$78$&$1$  \\
		 \hline
		&&$n=20$\\ 
		\hline 
\end{tabular}}
 \loigiai{
	\begin{itemchoice}
		\itemch \textbf{Đúng}.  \\
		 Khoảng biến thiên của mẫu số liệu là $80-60=20$.
		\itemch \textbf{Đúng}.\\ Số trung bình  cộng của mẫu số liệu được tính bởi công thức
		\[\overline{x}=\dfrac{8 \cdot 62+  9 \cdot 66+1 \cdot 70+1 \cdot 74+ 1\cdot 78}{20}=\dfrac{328}{5}.\]
		\itemch \textbf{Sai}. \\Phương sai của mẫu số liệu ghép nhóm đã  cho là
		 \[s^2=\dfrac{1}{20} \left(8 \cdot 62^2 + 9 \cdot 66^2+1 \cdot 70^2 +1 \cdot 74^2 +1 \cdot 78^2\right)-\left(\dfrac{328}{5}\right)^2=\dfrac{436}{25}.\]
		 \itemch \textbf{Sai}. \\
		Độ lệch chuẩn của mẫu số liệu ghép nhóm đã cho (làm tròn kết quả đến hàng phần mười  của kilôgam) là $s =\sqrt{\dfrac{436}{25}} \approx 4{,}2$ kg
	\end{itemchoice}
 }
\end{ex}

% nguồn: Lop_12_Dai_so_Chuong_3
\begin{ex}%[2D3H2-2]%[ma:1f1501f4]
	Thống kê tổng số giờ nắng trong tháng 9 tại một trạm quan trắc đặt ở Cà Mau trong các năm từ 2002 đến 2021 được thống kê như sau:\\ 
	 \centerline{
	\begin{tabular}{|*{10}{c|}}\hline
		$111{,}6$&  $134{,}9$& $130{,}3$& $134{,}2$& $140{,}9$& $109{,}3$& $152{,}4$& $156{,}3$& $116{,}1$& $96{,}7$\\\hline
		$105{,}2$& $80{,}8$& $80{,}8$&  $110$& $109$& $139$& $145$& $161$& $126$& $114$\\\hline
	\end{tabular}
	} 
	Người ta lập được  bảng tần số ghép nhóm như sau:\\
	 \centerline{
	\begin{tabular}{|*{6}{c|}}\hline
		Số giờ nắng& $[80;98)$&  $[98;116)$& $[116;134)$& $[134;152)$& $[152;170)$\\\hline
		Giá trị đại  diện& $89$& $107$& $125$& $143$& $161$\\
		Số năm& $3$& $6$& $3$&  $5$& $3$\\\hline
	\end{tabular}
	 }
	Xét  tính đúng sai của các khẳng định sau:
	\choiceTF 
		{\True Số trung bình của mẫu số liệu ghép nhóm  là $124{,}1$}
		{\True Phương sai của mẫu số liệu ghép nhóm là  $566{,}19$}
		{\True Độ lệch chuẩn của mẫu  số liệu ghép nhóm (kết quả các phép tính làm tròn đến hàng phần nghìn) là $23{,}795$}
		{\True Sai số tương đối của  độ lệch chuẩn của mẫu số liệu ghép nhóm so với độ lệch chuẩn của mẫu số liệu gốc (kết quả các phép tính làm tròn đến hàng phần nghìn) là $4{,}805$}
	 \loigiai{
		\begin{itemchoice}
			\itemch Đúng:  Xét theo mẫu số liệu gốc. Cỡ mẫu là $n=20$.\\
			Số trung bình của  mẫu số liệu trên là: $\overline{x}_1=\dfrac{111{,}6+134{,}9+\ldots +114}{20}=122{,}755$.\\
			Phương sai của mẫu số liệu trên là: $S_1^2=\dfrac{1}{20}(111{,}6^2+134{,}9^2+\ldots  +114^2)-122{,}755^2 \approx 515{,}453$.\\
			Độ lệch chuẩn của mẫu số liệu trên là:  $S_1 \approx \sqrt{515{,}453} \approx 22{,}704$.\\
			Xét theo mẫu số liệu  ghép nhóm\\
			Số trung bình của mẫu số liệu ghép nhóm là $\overline{x}_2=\dfrac{3\cdot 89+6\cdot 107+3\cdot 125+5\cdot 143+3\cdot  161}{20}=124{,}1$.
			\itemch Đúng: Phương sai của  mẫu số liệu ghép nhóm là\\
			$S_2^2=\dfrac{1}{20}(3\cdot  89^2+6\cdot 107^2+3\cdot 125^2+5\cdot 143^2+3\cdot 161^2)-124{,}1^2=566{,}19$.
			 \itemch Đúng: Độ lệch chuẩn của mẫu số liệu ghép nhóm là: $S_2=\sqrt{566{,}19} \approx 23{,}795$.
			\itemch Đúng: Sai số tương đối của độ lệch chuẩn của mẫu số liệu ghép nhóm so với độ lệch chuẩn của mẫu số liệu gốc là: $\dfrac{|S_2-S_1|}{S_1}=\dfrac{|23{,}795-22{,}704|}{22{,}704}\cdot 100 \approx 4{,}805$. 
		\end{itemchoice}
	 }
\end{ex}

% nguồn: Lop_12_Dai_so_Chuong_3
\begin{ex}%[2D3H2-2]%[ma:ea732d3a]
	Giá đóng cửa của một cổ phiếu là giá của cổ phiếu đó cuối một phiên giao dịch. Bảng sau thống kê giá đóng cửa (đơn vị: nghìn đồng) của hai mã cổ phiếu $A$ và $B$ trong 50 ngày giao  dịch liên tiếp.\\
	\centerline{ 
	\begin{tabular}{|*{6}{c|}}\hline
		Giá  đóng cửa& $[120;122)$& $[122;124)$& $[124;126)$& $[126;128)$& $[128;130)$\\\hline
		Cổ  phiếu A& $8$& $9$& $12$& $10$& $11$\\
		Cổ phiếu B& $16$& $4$& $3$&  $6$& $21$\\\hline
	\end{tabular}
	}\\ 
	Xét tính đúng sai của các khẳng định  sau:
	 \choiceTF
		{\True Xét mẫu số liệu của cổ phiếu $A$ ta có phương sai của mẫu  số liệu ghép nhóm là $7{,}5216$}
		{Xét mẫu số liệu của cổ phiếu $B$ ta có số trung bình của mẫu số liệu ghép nhóm là  $115{,}28$}
		{Xét mẫu số liệu của cổ phiếu $B$  ta có độ lệch chuẩn của mẫu số liệu ghép nhóm là $\sqrt{15{,}4096}$}
		{Người ta có thể dùng phương sai và độ lệch chuẩn để so sánh mức độ rủi ro của các loại cổ phiếu có giá trị trung bình  gần bằng nhau. Cổ phiếu nào có phương sai, độ lệch chuẩn cao hơn thì được coi là có độ rủi ro lớn hơn. Theo quan điểm trên, thì cổ phiếu $A$ có độ rủi ro thấp hơn cổ phiếu $B$}
	 \loigiai{
		Ta có bảng thống kê giá đóng  cửa theo giá trị đại diện:\\
		 \centerline{
		\begin{tabular}{|*{6}{c|}}\hline
			 Giá đóng cửa& $121$& $123$& $125$& $127$& $129$\\\hline
			 Cổ phiếu A& $8$& $9$& $12$& $10$& $11$\\
			 Cổ phiếu B& $16$& $4$& $3$& $6$& $21$\\\hline
		\end{tabular}
		} 
		\begin{itemchoice}
			\itemch Đúng: Xét mẫu số liệu của cổ  phiếu $A$:\\
			Số trung bình của  mẫu số liệu ghép nhóm là\\
			$\overline{x}_1=\dfrac{8\cdot 121+9\cdot 123+12\cdot  125+10\cdot 127+11\cdot 129}{50}=125{,}28$.\\
			Phương sai của  mẫu số liệu ghép nhóm là\\
			$S_1^2=\dfrac{1}{50}(8\cdot  121^2+9\cdot 123^2+12\cdot 125^2+10\cdot 127^2+11\cdot 129^2)-(125{,}28)^2=7{,}5216$.\\
			Độ lệch chuẩn của mẫu số liệu ghép nhóm  là $S_1=\sqrt{S_1^2}=\sqrt{7{,}5216}$.\\
			\itemch  Sai: Xét mẫu số liệu của cổ phiếu $B$:\\
			Số trung bình của  mẫu số liệu ghép nhóm là\\
			$\overline{x}_2=\dfrac{16\cdot 121+4\cdot 123+3\cdot 125+6\cdot 127+21\cdot  129}{50}=125{,}28.$\\
			 \itemch Sai: Phương sai của mẫu số liệu ghép nhóm là\\
			$S_2^2=\dfrac{1}{50}(16\cdot 121^2+4\cdot 123^2+3\cdot 125^2+6\cdot 127^2+21\cdot  129^2)-{{(125{,}48)}^2}=12{,}4096$.\\
			Độ lệch chuẩn của mẫu số liệu ghép nhóm là $S_2=\sqrt{S_2^2}=\sqrt{12{,}4096}$.\\ 
			\itemch Vậy nếu đánh giá độ rủi ro theo phương sai và độ lệch chuẩn thì  cổ phiếu $A$ có độ rủi ro thấp hơn cổ phiếu $B$.
		\end{itemchoice}
	 }
\end{ex}

% nguồn: Lop_12_Dai_so_Chuong_3
\begin{ex}%[2D3H2-2]%[ma:fbcca984]
	Thống  kê thời gian (đơn vị: phút) tập thể dục buổi sáng mỗi ngày trong tháng 4 năm 2024 của An cho kết quả như sau:
	\begin{center}
		\begin{tabular}{|>{\raggedright\arraybackslash}m{3.25cm}|>{\centering\arraybackslash}m{1.5cm}|>{\centering\arraybackslash}m{1.5cm}|>{\centering\arraybackslash}m{1.5cm}|>{\centering\arraybackslash}m{1.5cm}|>{\centering\arraybackslash}m{1.5cm}|}
			\hline 
			Thời gian (phút) & $[15;20)$ & $[20;25)$ & $[25;30)$ & $[30;35)$  & $[35;40)$ \\
			 \hline
			Số ngày & $5$ & $4$  & $10$ & $7$ & $4$ \\
			\hline 
		\end{tabular}
	\end{center}
	Khi đó 
	 \choiceTF
		{\True Khoảng biến thiên của mẫu số liệu về thời gian tập thể dục buổi sáng mỗi ngày  trong tháng 4 năm 2024 của bạn An là $25$}
		{Nhóm chứa tứ  phân vị thứ nhất là $[25;30)$}
		{\True Khoảng tứ phân vị của mẫu số liệu về thời gian tập thể dục buổi sáng mỗi ngày trong tháng 4 năm 2024  của bạn An là $9{,}375$}
		{Phương sai của mẫu số liệu là $36{,}14$ (kết quả được làm tròn  đến hàng phần trăm)}
	\loigiai{ 
		\begin{itemchoice}
			\itemch Đúng. Vì khoảng biến thiên của mẫu  số liệu về thời gian tập thể dục buổi sáng mỗi ngày trong tháng 4 năm 2024 của bạn An là $40-15=25$.
			\itemch  Sai. Vì
			\begin{itemize}
				\item Cỡ mẫu là: $30$. 
				\item Gọi $x_1;x_2;\ldots;x_30$ là mẫu số liệu gốc về thời gian tập thể dục buổi sáng mỗi  ngày trong tháng 4 năm 2024 của bạn An đã được sắp xếp theo thứ tự không giảm.
				\item Tứ phân vị thứ nhất của mẫu số liệu gốc là $x_8$ nên nhóm chứa tứ phân vị thứ nhất là  nhóm $[20;25)$.
			\end{itemize}
			 \itemch Đúng. Vì
			\begin{itemize}
				 \item Ta có $Q_1=20+\dfrac{\dfrac{1\cdot30}{4}-5}{4}\cdot5=23{,}125$
				\item  Tứ phân vị thứ ba của mẫu số liệu gốc là $x_{22}$ nên nhóm chứa tứ phân vị thứ ba là nhóm $[30;35)$ và ta có $Q_3=30+\dfrac{\dfrac{3\cdot30}{4}-19}{7}\cdot5=32{,}5$
				\item  Khoảng tứ phân vị của mẫu số liệu ghép nhóm là $\Delta_Q=Q_3-Q_1=32{,}5-23{,}125=9{,}375$.
			\end{itemize}
			\itemch Sai. Chọn giá trị đại diện cho các  nhóm số liệu, ta có
			\begin{center}
				\begin{tabular}{|>{\raggedright\arraybackslash}m{3.25cm}|>{\centering\arraybackslash}m{1.5cm}|>{\centering\arraybackslash}m{1.5cm}|>{\centering\arraybackslash}m{1.5cm}|>{\centering\arraybackslash}m{1.5cm}|>{\centering\arraybackslash}m{1.5cm}|}
					 \hline
					Thời gian (phút) & $[15;20)$ & $[20;25)$  & $[25;30)$ & $[30;35)$ & $[35;40)$ \\
					 \hline
					Giá trị đại diện & $17{,}5$ & $22{,}5$ & $27{,}5$  & $32{,}5$ & $37{,}5$\\
					\hline 
					Số ngày & $5$ & $4$ &  $10$ & $7$ & $4$ \\
					\hline 
				\end{tabular}
			\end{center}
			Thời gian trung bình tập thể dục buổi sáng mỗi  ngày của An là
			$$\overline{x}=\dfrac{1}{30}\cdot(17{,}5\cdot5+22{,}5\cdot4+27{,}5\cdot10+32{,}5\cdot7+37{,}5\cdot4)=\dfrac{83}{3}$$ 
			Phương sai của mẫu  số liệu là
			 $$s^2=\dfrac{1}{30}\cdot(17{,}5^2\cdot5+22{,}5^2\cdot4+27{,}5^2\cdot10+32{,}5^2\cdot7+37{,}5^2\cdot4)-\left(\dfrac{83}{3}\right)^2=\dfrac{1\,409}{36}\approx39{,}14$$
		\end{itemchoice}
	} 
\end{ex}


%%%%%%%%%%%%%%%% TRẢ LỜI NGẮN %%%%%%%%%%%%%%%%

% nguồn: Lop_12_Dai_so_Chuong_3
\begin{bt}%[2D3H2-2]%[ma:f69e5cf9]
	 \immini{
	Kết quả  khảo sát thời gian sử dụng liên tục (đơn vị: giờ) từ lúc sạc đầy cho đến khi hết của pin một số máy vi tính cùng loại được mô tả bằng biểu đồ bên.\\
	Hãy xác định độ lệch chuẩn của thời gian  sử dụng pin (kết quả được làm tròn đến hàng phần trăm)
	\shortans{0,19} 
	}{ 
	\begin{tikzpicture}[scale=.8,>=stealth,font=\footnotesize,line join=round,line cap=round,blue,x=10mm,y=7.5mm]
			\def\xmin{0} \def\xmax{6.5} 
			\def\ymin{0} \def\ymax{8} 
			\foreach \d in {1,2,...,7}{\draw[green,ultra thin]  (0,\d)--++(\xmax,0);}
			\draw[->] (\xmin,0)--(\xmax,0) node [below]{Thời  gian};
			\draw[->] (0,\ymin)--(0,\ymax) node [left]{Số  máy};
			\foreach \y in {0,1,...,7} 
			\draw (1pt,\y)--(-2pt,\y) node  [left] {$\y$};
			\fill[color=cyan,opacity=.75]  (1,0)--(1,2)--(2,2)--(2,4)--(3,4)--(3,7)--(4,7)--(4,5)--(5,5)--(5,0);
			\draw (1,0)--(1,2)--(2,2)--(2,4)--(3,4)--(3,7)--(4,7)--(4,5)--(5,5)--(5,0); 
			 \draw (2,0)--(2,2)
				 (3,0)--(3,4)
				 (4,0)--(4,5);
			\path (-1,\ymax+.75) node[right]{Thời  gian sử dụng pin của một số máy vi tính};
			\path (1,0)  node[below]{$7{,}2$}
				(2,0) node[below]{$7{,}4$} 
				(3,0) node[below]{$7{,}6$} 
				(4,0)  node[below]{$7{,}8$}
				(5,0)  node[below]{$8{,}0$}
				 (\xmax,0) node[below=4mm]{(giờ)}
				(0,\ymax-.5) node[left]{vi tính}; 
		\end{tikzpicture}
	} 
	\loigiai{ 
		 \centerline{
		\begin{tabular}{|c|c|c|c|c|}\hline
			 Giá trị đại diện& $7{,}3$& $7{,}5$& $7{,}7$& $7{,}9$\\\hline
			 Số máy& $2$& $4$& $7$& $5$\\\hline
		\end{tabular}
		}\\ 
		Cỡ  mẫu: ${{n}=18}$.\\
		Số trung bình: $\overline{x}=\dfrac{2\cdot 7{,}3+4\cdot 7{,}5+7\cdot 7{,}7+5\cdot  7{,}9}{18} \approx 7{,}67$.\\
		Phương sai: $S^2=\dfrac{2\cdot 7{,}3^2+4\cdot  7{,}5^2+7\cdot 7{,}7^2+5\cdot 7{,}9^2}{18}-7{,}67^2 \approx 0{,}04$.\\
	Độ lệch  chuẩn: $\sigma=\sqrt{0{,}04} \approx 0{,}19$.}
\end{bt}

% nguồn: Lop_12_Dai_so_Chuong_3
\begin{bt}%[2D3H2-2]%[ma:ff635b02]
	Bảng dưới đây thống kê  cự li ném tạ của một vận động viên.\\
	 \centerline{
	\begin{tabular}{|*{6}{c|}}\hline
		Cự li (m)& $[19;19{,}5)$& $[19{,}5;20)$&  $[20;20{,}5)$& $[20{,}5;21)$& $[21;21{,}5)$\\\hline
		Tần số& $13$& $45$& $24$& $12$& $6$\\\hline 
	\end{tabular}
	}\\ 
	Hãy tính độ lệch chuẩn của mẫu số liệu ghép nhóm trên (kết quả được làm tròn đến hàng  phần trăm)
	 \shortans{0,53}
	\loigiai{ 
		 \centerline{
		\begin{tabular}{|*{6}{c|}}\hline
			Giá trị đại diện& $19{,}25$& $19{,}75$& $20{,}25$& $20{,}75$& $21{,}25$\\\hline 
			Tần  số& $13$& $45$& $24$& $12$& $6$\\\hline
		\end{tabular}
		 }
		Cỡ mẫu:  ${{n}=100}$.\\
		Số trung bình: $\overline{x}=\dfrac{13\cdot 19{,}25+45\cdot 19{,}75+24\cdot  20{,}25+12\cdot 20{,}75+6\cdot 21{,}25}{100}=20{,}015$.\\
		Phương  sai: $S^2=\dfrac{13\cdot 19{,}25^2+45\cdot 19{,}75^2+24\cdot 20{,}25^2+12\cdot 20{,}75^2+6\cdot 21{,}25^2}{100}-20{,}015^2 \approx 0{,}28$.\\
		Độ  lệch chuẩn: ${\sigma=\sqrt{0{,}28} \approx 0{,}53}$.
	 }
\end{bt}

% nguồn: Lop_12_Dai_so_Chuong_3
\begin{bt}%[2D3H2-2]%[ma:22287c50]
	Bạn Mai ghi lại thời gian sử dụng điện  thoại di động mỗi ngày của mình trong $10$ ngày liên tiếp ở bảng sau (đơn vị: phút).
	\begin{center}
		\begin{tabular}{|l|l|l|l|l|l|l|l|l|l|}
			\hline $150$ & $251$ & $73$ & $188$ & $165$ & $225$ & $235$ & $144$ & $160$ & $244$  \\
			 \hline
		\end{tabular}
	\end{center}
	Bạn Mai ghép số  liệu trên thành $4$ nhóm có độ dài bằng nhau, với nhóm cuối cùng là $[220; 270$). Tính tỉ số giữa độ lệch chuẩn và trung bình mẫu của mẫu số liệu ghép nhóm (kết quả làm tròn đến hàng phần trăm).
	\shortans{0,29} 
	 \loigiai{
		Bảng số liệu ghép  nhóm:
		\begin{center}
			\begin{tabular}{|l|c|c|c|c|}
				\hline Thời gian sử dụng  & {$[70; 120)$} & {$[120; 170)$} & {$[170; 220)$} & {$[220; 270)$} \\
				\hline Số ngày & 1 & 4 & 1 & 4  \\
				 \hline
			\end{tabular}
		\end{center}
		Số trung  bình $\overline{x}=\dfrac{1\cdot 145+4\cdot 245+1\cdot 345+4\cdot 445}{10}=185$.\\
		Phương sai \[s^2=\dfrac{1\cdot (145-\overline{x})^2+4\cdot (245-\overline{x})^2+1\cdot (345-\overline{x})^2+4\cdot (445-\overline{x})^2}{10}=2\,900.\] 
		Độ lệch  chuẩn $s=\sqrt{s^2}=\sqrt{2\,900}$.\\
		Do đó $\dfrac{s}{\overline{x}}\approx 0{,}29$. 
	} 
\end{bt}

% nguồn: Lop_12_Dai_so_Chuong_3
\begin{bt}%[2D3H2-2]%[ma:39fb1c54]
	Một vận động viên luyện  tập chạy cự li 100 m đã ghi lại kết quả luyện tập như sau:\\
	 \centerline{
	\begin{tabular}{|c|c|c|c|c|}\hline
		Thời gian (giây)& $[10{,}2;10{,}4)$&  $[10{,}4;10{,}6)$& $[10{,}6;10{,}8)$& $[10{,}8;11)$\\\hline
		Số vận động  viên& $3$& $7$& $8$& $2$\\\hline
	\end{tabular}
	 }\\
	Hãy xác định phương sai của mẫu số liệu  trên?
	\shortans{0,03} 
	\loigiai{ 
		 Chọn giá trị đại diện cho các nhóm số liệu, ta có:\\
		 \centerline{
		\begin{tabular}{|c|c|c|c|c|}\hline
			Thời gian (giây)& $[10{,}2;10{,}4)$& $[10{,}4;10{,}6)$& $[10{,}6;10{,}8)$& $[10{,}8;11)$\\\hline 
			Giá trị đại diện&  $10{,}3$& $10{,}5$& $10{,}7$& $10{,}9$\\\hline
			Số vận  động viên& $3$& $7$& $8$& $2$\\\hline
		\end{tabular}
		} 
		Tổng số vận động  viên là: ${3+7+8+2=20}$\\
		Thời gian chạy trung bình của các vận động viên là:  \\
		$\overline{x}=\dfrac{1}{20}(10{,}3\cdot 3+10{,}5\cdot 7+10{,}7\cdot 8+10{,}9\cdot 2)=10{,}59$ (giây).\\ 
		Phương sai của  mẫu số liệu là:\\
	$s^2=\dfrac{1}{20}(10{,}3^2\cdot 3+10{,}5^2\cdot 7+10{,}7^2\cdot 8+10{,}9^2\cdot 2)-10{,}59^2=0{,}0299$.} 
\end{bt}

% nguồn: Lop_12_Dai_so_Chuong_3
\begin{bt}%[2D3H2-2]%[ma:307d71bc]
	\immini{Một siêu thị thống kê số tiền đơn vị: chục nghìn đồng) mà $44$ khách hàng mua hàng ở siêu thị đó trong một ngày. Số liệu được  ghi lại trong bảng sau. Phương sai của mẫu số liệu ghép nhóm trên là bằng bao nhiêu? (Làm tròn đến hàng phần mười).}{
	\begin{tabular}{|c|c|c|}
		 \hline
		Nhóm & Giá trị  đại diện & Tần số \\
		\hline 
		 $[40; 45)$ & $42{,}5$& $4$\\
		\hline 
		 $[45; 50)$ & $47{,}5$& $14$\\
		 \hline
		$[50; 55)$ &  $52{,}5$& $8$\\
		 \hline
		 $[55; 60)$ & $57{,}5$& $10$\\
		\hline 
		$[60; 65)$ & $62{,}5$& $6$\\ 
		\hline 
		$[65;  70)$ & $67{,}5$& $2$\\
		 \hline
		 & & $n=44$ \\
		 \hline
	\end{tabular}}
	\shortans{46,1} 
	 \loigiai{
		Số  trung bình cộng của mẫu số liệu ghép nhóm là
		\[\overline{x}=\dfrac{4\cdot 42{,}5+14\cdot 47{,}5+8\cdot 52{,}5+10\cdot 57{,}5+6\cdot  62{,}5+2\cdot 67{,}5}{44}=\dfrac{585}{11}.\]
		Vậy phương sai của của mẫu số liệu ghép  nhóm là
		\allowdisplaybreaks 
		\begin{eqnarray*}
			s^2&=&  \dfrac{4\cdot 42{,}5^2+14\cdot 47{,}5^2+8\cdot 52{,}5^2+10\cdot 57{,}5^2+6\cdot 62{,}5^2+2\cdot 67{,}5^2}{44}-\left(\dfrac{585}{11}\right)^2\\
			&=&\dfrac{22\,325}{484}\approx46{,}1. 
		\end{eqnarray*}
	 }
\end{bt}

% nguồn: Lop_12_Dai_so_Chuong_3
\begin{bt}%[2D3H2-2]%[ma:fae26e6e]
	Hãy tính độ lệch chuẩn của mẫu số liệu ghép nhóm (kết  quả được làm tròn đến hàng phần trăm).
	\begin{center}
		\begin{tabular}{|c|c|c|c|c|c|}
			\hline Chiều  cao $(\mathrm{cm})$ & {$[160; 164)$} & {$[164; 168)$} & {$[168; 172)$} & {$[172; 176)$} & {$[176; 180)$} \\
			\hline Số học $\sinh$ & $3$ & $5$ & $8$ & $4$  & $1$ \\
			 \hline
		\end{tabular}
	\end{center}
	 \shortans{4,26}
	 \loigiai{
		Cỡ mẫu $n=21$.\\ 
		Giá trị trung bình của mẫu  số liệu mới
		\[\overline{x}=\dfrac{1}{n}\left(n_1c_1+n_2c_2+\ldots  +n_kc_k\right)=\dfrac{1}{21}\left(3\cdot 162+5\cdot 166+8\cdot 170+4\cdot 174+1\cdot 178\right)=\dfrac{3550}{21}.\]
		 Phương sai của mẫu số liệu mới
		\begin{eqnarray*}
			&S^2&=\dfrac{1}{n}\left[n_1\left(c_1-\bar{x}\right)^2+n_2\left(c_2-\bar{x}\right)^2+\ldots +n_k\left(c_k-\bar{x}\right)^2\right]\\ 
			&&=\dfrac{1}{21}\left[3\left(162-\dfrac{3550}{21}\right)^2+5\left(166-\dfrac{3550}{21}\right)^2+\cdots +1\left(178-\dfrac{3550}{21}\right)^2\right]\\ 
			&&=\dfrac{8000}{441}. 
		\end{eqnarray*}
		Độ lệch chuẩn của mẫu số liệu  mới $\sigma =\sqrt{S^2}=\sqrt{\dfrac{8000}{441}}=\dfrac{40\sqrt{5}}{21}\approx 4,26$.
	} 
\end{bt}

% nguồn: Lop_12_Dai_so_Chuong_3
\begin{bt}%[2D3H2-2]%[ma:9d50cc61]
	 Thành tích môn nhảy cao của các vận động viên tại một giải điền kinh dành cho học sinh trung học phổ thông như sau:\\
	 \centerline{
	\begin{tabular}{|c|c|c|c|c|}\hline
		Mức xà  (cm)& $[170;172)$& $[172;174)$& $[174;176)$& $[176;178)$\\\hline
		 Số vận động viên& $3$& $10$& $6$& $1$\\\hline
	\end{tabular}
	 }
	Hãy xác định độ lệch chuẩn của thời gian sử dụng pin (kết  quả được làm tròn đến hàng phần trăm)
	 \shortans{1,53}
	 \loigiai{
		 Mẫu số liệu với giá trị đại diện\\
		\centerline{ 
		\begin{tabular}{|c|c|c|c|c|}\hline
			Mức xà (cm)& $[170;172)$& $[172;174)$&  $[174;176)$& $[176;178)$\\\hline
			Giá trị đại  diện& $171$& $173$& $175$& $177$\\\hline
			 Số vận động viên& $3$& $10$& $6$& $1$\\\hline
		\end{tabular}
		 }\\
		Giá  trị trung bình: $\overline{x}=\dfrac{1}{20}(171\cdot 3+173\cdot 10+175\cdot 6+177\cdot 1)=173{,}5$ (cm).\\
		Phương sai của mẫu  số liệu: $s^2=\dfrac{1}{20}(171^2\cdot 3+173^2\cdot 10+175^2\cdot 6+177^2\cdot 1)-173{,}5^2=2{,}35$.\\
	Độ lệch chuẩn của mẫu số liệu: $s=\sqrt{2{,}35}=\dfrac{\sqrt{235}}{10} \approx 1{,}53$ (cm).} 
\end{bt}

% nguồn: Lop_12_Dai_so_Chuong_3
\begin{bt}%[2D3H2-2]%[ma:5ad26b59]
	Bác Hải đo chiều cao của $50$ cây giống được lựa chọn ngẫu nhiên. Kết quả được ghi lại  trong bảng số liệu ghép nhóm dưới đây.
	\begin{center}
		\begin{tabular}{|c|c|c|c|c|c|}
			\hline Chiều cao cm & {$[15; 16)$} & {$[16; 17)$}  & {$[17; 18)$} & {$[18; 19)$} & {$[19; 20)$} \\
			\hline Tần số tưong đối & $12 \%$ & $24 \%$ & $36 \%$ & $20 \%$ & $8 \%$ \\ 
			\hline 
		\end{tabular}
	\end{center}
	Tính tỉ số của độ lệch chuẩn và số trung bình của mẫu số liệu  ghép nhóm trên (kết quả làm tròn đến hàng phần trăm).
	 \shortans{0,06}
	\loigiai{ 
		Bảng tần số  ghép nhóm
		\begin{center}
			\begin{tabular}{|c|c|c|c|c|c|}
				\hline 
				Nhóm & {$[15;16)$} & {$[16; 17)$} & {$[17; 18)$} &  {$[18; 19)$} & {$[19; 20)$} \\
				 \hline
				Tần số &  $6$ & $12$ & $18$ & $10$ & $4$ \\
				\hline 
			\end{tabular}
		\end{center}
		Số trung bình của  mẫu số liệu ghép nhóm là $$\overline{x}=\dfrac{15{,}5\cdot6+16{,}5\cdot12+17{,}5\cdot18+18{,}5\cdot10+19{,}5\cdot4}{50}=17{,}38.$$\\
		Độ lệch chuẩn  của mẫu số liệu ghép nhóm là
		\begin{eqnarray*}
			S&=&\dfrac{1}{50}[6\cdot(15{,}5-17{,}38)^2+12\cdot(16{,}5-17{,}38)^2+18\cdot(17{,}5-17{,}38)^2\\ 
			&&+10\cdot(18{,}5-17{,}38)^2+4\cdot(19{,}5-17{,}38)^2]\\ 
			 S&\approx &1{,}107068.
		\end{eqnarray*}
	Tỉ số của độ lệch chuẩn và số trung bình của mẫu số liệu ghép nhóm trên là $$\dfrac{S}{\overline{x}} \approx \dfrac{1{,}107068}{17{,}38} \approx  0{,}06.$$}
\end{bt}

% nguồn: Lop_12_Dai_so_Chuong_3
\begin{bt}%[2D3H2-2]%[ma:2e2a36f6]
	Để đánh giá chất lượng một loại pin điện thoại mới, người ta ghi lại thời gian nghe nhạc liên  tục của điện thoại được sạc đầy pin cho đến khi hết pin cho kết quả như sau:\\
	\centerline{ 
	\begin{tabular}{|*{6}{c|}}\hline
		Thời gian (giờ)& $[5;5{,}5)$& $[5{,}5;6)$& $[6;6{,}5)$&  $[6{,}5;7)$& $[7;7{,}5)$\\\hline
		Số chiếc  điện thoại
		(tần số)& $2$& $8$& $15$& $10$&  $5$\\\hline
	\end{tabular}
	}\\ 
	Tính độ lệch chuẩn của mẫu số liệu  ghép nhóm trên (kết quả được làm tròn đến hàng phần trăm)
	\shortans{0,53} 
	\loigiai{ 
		Chọn  giá trị đại diện cho các nhóm số liệu, ta có:\\
		\centerline{ 
		\begin{tabular}{|*{6}{c|}}\hline
			Thời gian (giờ)& $[5;5{,}5)$& $[5{,}5;6)$& $[6;6{,}5)$& $[6{,}5;7)$&  $[7;7{,}5)$\\\hline
			Giá trị đại diện& $5{,}25$& $5{,}75$& $6{,}25$&  $6{,}75$& $7{,}25$\\\hline
			Số  chiếc điện thoại
			(tần số)& $2$&  $8$& $15$& $10$& $5$\\\hline
		\end{tabular}
		}\\ 
		Thời gian trung bình nghe  nhạc liên tục của điện thoại là: \\
		$\overline{x}=\dfrac{1}{40}(5{,}25\cdot 2+5{,}75\cdot 8+6{,}25\cdot 15+6{,}75\cdot 10+7{,}25\cdot  5)=6{,}35.$\\
		Phương sai  của mẫu số liệu là:\\
		$s^2=\dfrac{1}{40}(5{,}25^2\cdot 2+5{,}75^2\cdot  8+6{,}25^2\cdot 15+6{,}75^2\cdot 10+7{,}25^2\cdot 5)-6{,}35^2=0{,}2775.$\\
	Độ lệch chuẩn  của mẫu số liệu là: $\sqrt{0{,}2775}=\dfrac{\sqrt{111}}{20} \approx 0{,}53$.}
\end{bt}

% nguồn: Lop_12_Dai_so_Chuong_3
\begin{bt}%[2D3H2-2]%[ma:024a1bd3]
	Một bác tài xế thống kê lại  độ dài quãng đường (đơn vị: km) bác đã lái xe mỗi ngày trong một tháng ở bảng sau:\\
	 \centerline{
	\begin{tabular}{|*{6}{c|}}\hline
		Độ dài quãng  đường (km)& $[50;100)$& $[100;150)$& $[150;200)$& $[200;250)$& $[250;300)$\\\hline
		Số ngày& $5$& $10$&  $9$& $9$& $2$\\\hline
	\end{tabular}
	}\\ 
	Hãy xác định độ lệch  chuẩn của mẫu số liệu trên
	\shortans{3100} 
	\loigiai{ 
		\centerline{ 
		\begin{tabular}{|*{6}{c|}}\hline
			Độ dài quãng đường (km)& $75$& $125$& $175$& $225$& $275$\\\hline 
			Số ngày& $5$& $10$& $9$& $9$& $2$\\\hline 
		\end{tabular}
		}\\ 
		Số trung  bình: $\overline{x}=\dfrac{5\cdot 75+10\cdot 125+9\cdot 175+4\cdot 225+2\cdot 275}{30}=155$.\\
	Độ lệch  chuẩn: $\sigma =\sqrt{\dfrac{5\cdot 75^2+10\cdot 125^2+9\cdot 175^2+4\cdot 225^2+2\cdot 275^2}{30}-155^2}=3100$.}
\end{bt}

% nguồn: Lop_12_Dai_so_Chuong_3
\begin{bt}%[2D3H2-2]%[ma:127679eb]
	Cân nặng của một số quả mít trong một khu vườn được thống  kê ở bảng sau:\\
	\centerline{ 
	\begin{tabular}{|*{6}{c|}}\hline
		Cân nặng (kg)& $[4;6)$& $[6;8)$& $[8;10)$& $[10;12)$& $[12;14)$\\\hline 
		Số quả mít& $6$& $12$& $19$& $9$& $4$\\\hline 
	\end{tabular}
	}\\ 
	Hãy tính  phương sai của mẫu số liệu ghép nhóm trên. (Kết quả các phép tính làm tròn đến hàng phần trăm)?
	\shortans{4,80} 
	 \loigiai{
		Ta có bảng thống kê cân nặng của các quả mít theo giá trị đại diện:\\ 
		\centerline{ 
		\begin{tabular}{|*{6}{c|}}\hline
			Cân  nặng đại diện (kg)& $5$& $7$& $9$& $11$& $13$\\\hline
			 Tần số& $6$& $12$& $19$& $9$& $4$\\\hline
		\end{tabular}
		}\\ 
		Cỡ mẫu  ${n=6+12+19+9+4=50}$.\\
		Số trung bình của mẫu số liệu ghép nhóm là: $\overline{x}=\dfrac{6\cdot 5+12\cdot 7+19\cdot 9+9\cdot 11+4\cdot  13}{50}=8{,}72$.\\
		Phương sai của mẫu số liệu  ghép nhóm là: \\
		$S^2=\dfrac{1}{50}(6\cdot 5^2+12\cdot 7^2+19\cdot 9^2+9\cdot 11^2+4\cdot  13^2)-8{,}72^2\approx 4{,}80$.
	} 
\end{bt}

% nguồn: Lop_12_Dai_so_Chuong_3
\begin{bt}%[2D3H2-2]%[ma:d8d60e34]
	Hãy tính độ lệch chuẩn của mẫu số liệu  ghép nhóm (kết quả được làm tròn đến hàng phần trăm) \\
	\centerline{ 
	\begin{tabular}{|*{6}{c|}}\hline
		Chiều cao (cm)&  $[160;164)$& $[164;168)$& $[168;172)$& $[172;176)$& $[176;180)$\\\hline
		Số học sinh& $3$& $5$& $8$&  $4$& $1$\\\hline
	\end{tabular}
	 }
	 \shortans{4,26}
	 \loigiai{
		Cỡ mẫu: ${n=21}$. Giá trị  trung bình của mẫu số liệu mới:\\
		$\overline{x}=\dfrac{1}{n}(n_1c_1+n_2c_2+\ldots +n_kc_k)=\dfrac{1}{21}(3\cdot 162+5\cdot 166+8\cdot 170+4\cdot 174+1\cdot 178) =\dfrac{3550}{21}$.\\ 
		Phương sai của mẫu  số liệu mới: $S^2=\dfrac{1}{n}\left[n_1{{\left(c_1-\overline{x}\right)}^2}+n_2{{\left(c_2-\overline{x}\right)}^2}+\ldots +n_k{{\left(c_k-\overline{x}\right)}^2}\right]$.\\
		$=\dfrac{1}{21}\left[3{{\left(162-\dfrac{3550}{21}\right)}^2}+5{{\left(166-\dfrac{3550}{21}\right)}^2}+\ldots +1{{\left(178-\dfrac{3550}{21}\right)}^2}\right]=\dfrac{8000}{441}$.\\ 
	Độ lệch chuẩn của mẫu số liệu mới: $\sigma  =\sqrt{S^2}=\sqrt{\dfrac{8000}{441}}=\dfrac{40\sqrt{5}}{21}\approx 4{,}26$.}
\end{bt}

% nguồn: Lop_12_Dai_so_Chuong_3
\begin{bt}%[2D3H2-2]%[ma:3d404ab4]
	Bảng bên dưới cho ta bảng tần số ghép nhóm về số liệu thống kê chiều cao (đơn vị: mét) của $40$ núi cao nhất Đông Nam  Á.\\
	 \centerline{\textit{(Nguồn: \href{https://vi.wikipedia.ong}{https://vi.wikipedia.ong})}}
	\begin{center}
		\begin{tabular}{|c|c|}
			 \hline
			Nhóm & Tần số  \\
			 \hline
			$[3\,500; 4\,000)$ & $10$  \\
			$[4\,000; 4\,500)$ & $7$  \\
			$[4\,500;  5\,000)$ & $16$ \\
			$[5\,000;  5\,500)$ & $4$ \\
			$[5\,500; 6\,000)$ & $3$ \\ 
			\hline 
			& $n=40$  \\
			\hline 
		\end{tabular}
	\end{center}
	 Độ lệch chuẩn của mẫu số liệu đó bằng bao nhiêu (làm tròn kết quả đến hàng đơn vị)?
	\shortans{590} 
	\loigiai{ 
		 Ta có bảng sau
		\begin{center}
			\begin{tabular}{|c|c|c|}
				 \hline
				Nhóm & Giá  trị đại diện & Tần số \\
				\hline 
				$[3\,500; 4\,000)$& $3\,750$ &  $10$ \\
				$[4\,000; 4\,500)$& $4\,250$ & $7$ \\ 
				$[4\,500; 5\,000)$&  $4\,750$ & $16$ \\
				$[5\,000; 5\,500)$& $5\,250$  & $4$ \\
				$[5\,500; 6\,000)$ & $5\,750$&  $3$ \\
				 \hline
			\end{tabular}
		\end{center}
		Suy ra $\overline{x}= \dfrac{10 \cdot 3\,750 + 7 \cdot 4\,250 + 16 \cdot 4\,750 + 4 \cdot 5\,250 + 3 \cdot 5\, 750 }{40}= 4\,537{,}  5$.\\
		Suy ra độ lệch chuẩn của mẫu số liệu là \[s=\sqrt{\dfrac{10 \cdot 3\,750^2 + 7 \cdot 4\,250^2 + 16 \cdot 4\,750^2 + 4 \cdot 5\,250^2 + 3 \cdot 5\, 750^2 }{40} -(4\,537{,}5)^2}\approx  590 \,(\text{m}).\]
	 }
\end{bt}

% nguồn: Lop_12_Dai_so_Chuong_3
\begin{bt}%[2D3H2-2]%[ma:bb567501]
	Chiều dài của 40 bé trai sơ sinh 12 ngày tuổi  chọn ngẫu nhiên ở một bệnh viện được nhà nghiên cứu thống kê trong bảng dưới đây:\\
	\centerline{ 
	\begin{tabular}{|*{7}{c|}}\hline
		Chiều dài (cm)& $[44;46)$& $[46;48)$& $[48;50)$& $[50;52)$& $[52;54)$& $[54;56)$\\\hline 
		Số trẻ& $3$& $3$& $10$& $15$& $7$&  $2$\\\hline
	\end{tabular}
	}\\ 
	Tính độ lệch chuẩn của chiều dài nhóm 40 bé  trai sơ sinh (làm tròn kết quả đến hàng phần nghìn)
	\shortans{2,43} 
	\loigiai{ 
		Bổ sung thêm các giá trị đại diện vào bảng trên ta lập được bảng sau. Chiều dài của 40  bé trai sơ sinh (đơn vị: cm)\\
		\centerline{ 
		\begin{tabular}{|c|c|c|}\hline
			 Nhóm& $c_i$& $n_i$\\\hline
			$[44;46)$& $45$& $3$\\ 
			$[46;48)$& $47$&  $3$\\
			$[48;50)$& $49$& $10$\\ 
			$[50;52)$& $51$&  $15$\\
			$[52;54)$&  $53$& $7$\\
			$[54;56)$& $55$& $2$\\\hline 
			&&$n=40$\\\hline 
		\end{tabular}
		} 
		Từ mẫu số  liệu đã cho, ta tính được số trung bình là:\\
		$\overline{x}=\dfrac{3\cdot 45+3\cdot  47+10\cdot 49+15\cdot 51+7\cdot 53+2\cdot 55}{40}=\dfrac{2012}{40}=50{,}3$.\\
		$\overline{x}$ không phải là  số nguyên nên để tính phương sai ta tính:\\
		$\overline{x^2}=\dfrac{3\cdot 45^2+3\cdot 47^2+10\cdot 49^2+15\cdot 51^2+7\cdot  53^2+2\cdot 55^2}{40}=2536.$\\
		 Do đó:${s^2=\overline{x^2}-(\overline{x})^2 =2536-50{,}3^2=2536-2530{,}09 =5{,}91}$.\\
	Vậy mẫu  số liệu về chiều dài của 40 trẻ sơ sinh có độ lệch chuẩn là ${s=\sqrt{5{,}91} \approx 2{,}43}$}
\end{bt}

% nguồn: Lop_12_Dai_so_Chuong_3
\begin{bt}%[2D3V2-2]%[ma:daf1c5af]
	Bảng dưới đây thống kê cự li ném tạ của một  vận động viên.
	\begin{center}
		\begin{tabular}{|c|c|c|c|c|c|}
			 \hline
			Cự li (m)  & $[19; 19{,}5)$ & $[19{,}5; 20)$ & $[20; 20{,}5)$ & $[20{,}5; 21)$ & $[21; 21{,}5)$ \\
			 \hline
			Tần số & $13$ & $45$ & $24$ & $12$ & $6$ \\ 
			\hline 
		\end{tabular}
	\end{center}
	Phương sai của mẫu số liệu  ghép nhóm trên bằng bao~nhiêu? (Kết quả làm tròn đến hàng phần trăm)
	 \shortans{0,28}
	 \loigiai{
		Cỡ mẫu $n=100$  \\
		Số trung bình: $\bar{x}=\dfrac{13\cdot 19{,}25+45\cdot 19{,}75+24\cdot 20{,}25+12\cdot 20{,}75+6\cdot 21{,}25}{100}=20{,}015$. \\ 
		Phương sai: $S^2=\dfrac{13\cdot 19{,}25^2+45\cdot 19{,}75^2+24\cdot 20{,}25^2+12\cdot  20{,}75^2+6\cdot 21{,}25^2}{100}-20{,}015^2\approx 0{,}28$
	 }
\end{bt}


%%%%%%%%%%%%%%%% TỰ LUẬN %%%%%%%%%%%%%%%%

% nguồn: Lop_12_Dai_so_Chuong_3
\begin{bt}%[2D3H2-2]%[ma:7773a316]
	Cân nặng của  một số quả mít trong một khu vườn được thống kê ở bảng sau
	\begin{longtable}{|c|c|c|c|c|c|}
		 \hline
		Cân  nặng & $[4;6)$ & $[6;8)$ & $[8;10)$ & $[10;12)$ & $[12;14)$\\
		 \hline
		Số quả  mít & $6$ & $12$ & $19$ & $9$ & $4$\\
		\hline 
	\end{longtable}
	\noindent Hãy tính phương sai và độ  lệch chuẩn của mẫu số liệu ghép nhóm trên. (Kết quả các phép tính làm tròn đến hàng phần trăm.)
	\loigiai{ 
		Ta có bảng thống kê cân nặng của các quả mít theo giá trị  đại diện
		\begin{longtable}{|c|c|c|c|c|c|}
			\hline 
			Câng nặng đại diện (kg) & $5$ & $7$ & $9$ & $11$ &  $13$\\
			\hline 
			Số quả mít & $6$ & $12$ & $19$ & $9$ &  $4$\\
			 \hline
		\end{longtable}
		\noindent Cỡ  mẫu $ n=6+12+19+9+4=50$.\\
		 Số trung bình của mẫu số liệu ghép nhóm là $\overline{x}=\dfrac{6\cdot 5+12\cdot 7+19\cdot 9+9\cdot 11+4\cdot 13}{50}=8{,}72.$\\
		Phương sai của  mẫu số liệu ghép nhóm là\\
		$$S^2=\dfrac{1}{50}\left(6\cdot{5^2}+12\cdot{7^2}+19\cdot 9^2+9\cdot{11^2}+4\cdot 13^2\right)-8{,}72^2\approx  4{,}80.$$
		Độ lệch chuẩn của mẫu số liệu ghép nhóm là $ S\approx\sqrt{4{,}80}\approx  2{,}19.$
	 }
\end{bt}

% nguồn: Lop_12_Dai_so_Chuong_3
\begin{bt}%[2D3H2-2]%[ma:ff38f0eb]
	Bảng dưới đây thống kê cự li ném tạ của  một vận động viên.
	\begin{longtable}{|c|c|c|c|c|c|}
		 \hline
		 Cự li (m) & $[19;19{,}5)$ & $[19{,}5;20)$ & $[20;20{,}5)$ & $[20{,}5;21)$ & $[21;21{,}5)$\\
		\hline 
		Tần số & $13$ & $45$ & $24$ & $12$ &  $6$\\
		\hline 
	\end{longtable}
	Hãy tính phương sai  và độ lệch chuẩn của mẫu số liệu ghép nhóm trên.
	 \loigiai{
		\begin{longtable}{|c|c|c|c|c|c|}
			\hline 
			Giá trị  đại diện & $19{,}25$ & $19{,}75$ & $20{,}25$ & $20{,}75$ & $21{,}25$\\
			\hline 
			Tần số & $13$ & $45$  & $24$ & $12$ & $6$\\
			 \hline
		\end{longtable}
		Cỡ  mẫu $ n=100$.\\
		Số trung bình $\overline{x}=\dfrac{13\cdot 19{,}25+45\cdot 19{,}75+24\cdot 20{,}25+12\cdot 20{,}75+6\cdot 21{,}25}{100}=20{,}015$.\\ 
		Phương sai $S^2=\dfrac{13\cdot 19{,}25^2+45\cdot 19{,}75^2+24\cdot 20{,}25^2+12\cdot 20{,}75^2+6\cdot  21{,}25^2}{100}-20{,}015^2\approx 0{,}28$.\\
		Độ lệch chuẩn  $s=\sqrt{0{,}28}\approx 0{,}53$.\\
	} 
\end{bt}

% nguồn: Lop_12_Dai_so_Chuong_3
\begin{bt}%[2D3H2-2]%[ma:8290bc57]
	Thống kê tổng số giờ nắng trong tháng 9 tại một trạm quan  trắc đặt ở Cà Mau trong các năm từ 2002 đến 2021 được thống kê như sau
	\begin{longtable}{|c|c|c|c|c|c|c|c|c|c|}
		\hline 
		 $111{,}6$ & $134{,}9$ & $130{,}3$ & $134{,}2$ & $140{,}9$ & $109{,}3$ & $154{,}4$ & $156{,}3$ & $116{,}1$ & $96{,}7$ \\
		\hline 
		$105{,}2$ & $80{,}8$ & $80{,}8$ & $110$ &  $109$ & $139$ & $145$ & $161$ & $126$ & $114$\\
		 \hline
	\end{longtable}
	\begin{enumerate}
		\item Hãy tính phương sai và độ lệch chuẩn của mẫu số liệu trên. 
		\item Hãy lập bảng tần  số ghép nhóm với nhóm đầu tiên là $\left[80;98\right)$ và độ dài mỗi nhóm bằng $18$. Tính phương sai, độ lệch chuẩn của mẫu số liệu ghép nhóm.
		\item Hãy tính sai số tương đối của độ lệch chuẩn của mẫu số liệu ghép nhóm so với độ lệch chuẩn của mẫu số liệu gốc. (Kết quả các phép  tính làm tròn đến hàng phần nghìn.)
	\end{enumerate}
	 \loigiai{
		\begin{enumerate}
			\item Cỡ mẫu là $  n=20$.\\
			Số trung bình của mẫu  số liệu trên là $\overline{x}_1=\dfrac{111{,}6+134{,}9+\ldots+114}{20}=122{,}755.$\\
			Phương sai  của mẫu số liệu trên là
			$$S_1^2=\dfrac{1}{20}\left(111{,}6^2+134{,}9^2+\ldots+114^2\right)-122{,}755^2\approx  515{,}453.$$
			Độ  lệch chuẩn của mẫu số liệu trên là $S_1\approx\sqrt{515{,}453}\approx 22{,}704.$
			 \item Ta có bảng sau
			\begin{longtable}{|c|c|c|c|c|c|}
				\hline 
				Số giờ nắng & $[80;98)$ & $[98;116)$ &  $[116;134)$ & $[134;152)$ & $[152;170)$\\
				 \hline
				Giá trị đại  diện & $89$ & $107$ & $125$ & $143$ & $161$\\
				 \hline
				Số năm & $3$ & $6$  & $3$ & $5$ & $3$\\
				\hline 
			\end{longtable}
			 Số trung bình của mẫu số liệu ghép nhóm là $$\overline{x}_2=\dfrac{3\cdot 89+6\cdot 107+3\cdot 125+5\cdot 143+3\cdot 161}{20}=124{,}1.$$
			Phương sai của mẫu  số liệu ghép nhóm là
			$$S_2^2=\dfrac{1}{20}\left(3\cdot 89^2+6\cdot 107^2+3\cdot  125^2+5\cdot 143^2+3\cdot 161^2\right)-124{,}1^2=566{,}19.$$
			Độ lệch  chuẩn của mẫu số liệu ghép nhóm là $S_2=\sqrt{566{,}19}\approx 23{,}795$.
			\item Sai số tương đối của độ lệch chuẩn của mẫu số liệu ghép nhóm  so với độ lệch chuẩn của mẫu số liệu gốc là $\dfrac{\left|S_2-S_1\right|}{S_1}=\dfrac{\left| 23{,}795-22{,}704\right|}{22{,}704}\cdot 100\%\approx 4{,}805\%$.
		\end{enumerate}
	} 
\end{bt}
